\documentclass[10pt,a4paper]{article}
\usepackage[utf8]{inputenc}
\usepackage[T1]{fontenc}
\usepackage[spanish,es-nodecimaldot]{babel}
\usepackage{lmodern}
\usepackage[margin=1.55cm]{geometry}
\usepackage{enumitem}
\usepackage{tabularx}
\usepackage{array}
\usepackage{xcolor}
\usepackage{titlesec}
\usepackage{hyperref}
\usepackage{microtype}
\usepackage{fancyhdr}
\usepackage{lastpage}
\usepackage{ragged2e}

\definecolor{accent}{HTML}{2A5DB0}
\definecolor{textgray}{HTML}{444444}
\hypersetup{colorlinks=true,urlcolor=accent,linkcolor=accent,pdfauthor={Gabriel Andrés Salgado Maldonado},pdftitle={Currículum Vitae - Gabriel Salgado Maldonado}}
\setlength{\parindent}{0pt}
\setlength{\parskip}{2pt}
\setlist[itemize]{leftmargin=1.15em,itemsep=1.5pt,topsep=2pt,parsep=0pt}
\titleformat{\section}{\large\bfseries\color{accent}}{}{0pt}{}[\vspace{-3pt}\color{accent}\titlerule]
\titlespacing*{\section}{0pt}{10pt}{5pt}
\pagestyle{fancy}
\fancyhf{}
\renewcommand{\headrulewidth}{0pt}
\fancyfoot[C]{\scriptsize Gabriel Salgado Maldonado \hfill Página \thepage\ de \pageref{LastPage}}

\newcommand{\entry}[4]{%
\textbf{#1}\hfill {\small #2}\\
{\emph{#3}}\\[-1pt]
#4\par\vspace{4pt}}

\newcommand{\pub}[2]{\hangindent=1.1em\hangafter=1 #1\\[-1pt]{\small\color{textgray}#2}\par\vspace{4pt}}

\begin{document}

{\Huge\bfseries Gabriel Andrés Salgado Maldonado}\\[3pt]
{\large\color{textgray}Fonoaudiólogo · Epidemiólogo · Académico}\\[6pt]
Santiago, Chile \quad | \quad \href{mailto:gsalgadom@uandes.cl}{gsalgadom@uandes.cl} \quad | \quad \href{mailto:gsalgadoma@gmail.com}{gsalgadoma@gmail.com} \quad | \quad +56 9 6237 3080\\
\href{https://orcid.org/0000-0002-4897-8863}{ORCID: 0000-0002-4897-8863} \quad | \quad \href{https://scholar.google.com/citations?user=lUJ6RBoAAAAJ}{Google Scholar} \quad | \quad \href{https://gsalgadoma.github.io}{gsalgadoma.github.io}

\section*{Perfil profesional}
Fonoaudiólogo y Magíster en Epidemiología con trayectoria clínica, académica y de gestión en salud. Experiencia en deglución, comunicación, voz, vía aérea, neurología cognitiva y rehabilitación en contextos hospitalarios y de alta complejidad. Desarrollo académico en docencia de pregrado, posgrado y educación continua, investigación aplicada, salud pública y liderazgo profesional.

\section*{Educación}
\entry{Magíster en Epidemiología}{2017--2018}{Universidad de los Andes, Facultad de Medicina, Departamento de Salud Pública y Epidemiología · Santiago, Chile}{}
\entry{Fonoaudiólogo, titulado con distinción}{2009}{Universidad del Desarrollo, Escuela de Fonoaudiología · Concepción, Chile}{}
\entry{Licenciado en Fonoaudiología}{2008}{Universidad del Desarrollo, Escuela de Fonoaudiología · Concepción, Chile}{}

\section*{Postítulos y formación de posgrado}
\begin{itemize}
  \item \textbf{Diplomado en Vocología}, Universidad de los Andes, 2025.
  \item \textbf{Diplomado en Docencia Universitaria}, Universidad de los Andes, 2024.
  \item \textbf{Diplomado de Unidad de Cuidados Intensivos del Adulto}, AST-Educación, 220 h, 2022.
  \item \textbf{Diplomado de Gestión y Calidad en Salud}, AST-Educación, 200 h, 2022.
  \item \textbf{Trastornos Deglutorios y Disfagia: Diagnóstico y Tratamiento Interdisciplinario}, Fundación Fono, Argentina, 320 h, 2012.
  \item \textbf{Diplomado en Neuropsicología y Neuropsiquiatría del Adulto}, Pontificia Universidad Católica de Chile / Universidad de Chile, 100 h, 2011.
\end{itemize}

\section*{Posiciones académicas actuales}
\entry{Profesor Asistente Extraordinario}{nov. 2025--presente}{Universidad de los Andes, Facultad de Medicina, Escuela de Fonoaudiología}{Docencia en Salud Pública, Rehabilitación Basada en la Comunidad y posgrado.}
\entry{Docente Instructor}{ago. 2025--presente}{Universidad de los Andes, Facultad de Enfermería y Obstetricia, Escuela de Enfermería}{Coordinación del Minor de Salud Pública y docencia en Epidemiología y Salud Pública, Fisiología, Seminario de Investigación y Bioestadística.}
\entry{Director, Diplomado de Deglutología del Adulto}{may. 2026--presente}{Universidad de los Andes, Escuela de Fonoaudiología}{Dirección académica del programa, coordinación de actividades y docencia.}
\entry{Director, Diplomado de Fonoaudiología en Medicina Intensiva}{may. 2024--presente}{Universidad de los Andes, Escuela de Fonoaudiología}{Dirección académica, coordinación y docencia de posgrado.}

\section*{Experiencia profesional}
\entry{Fonoaudiólogo}{abr. 2023--abr. 2026}{Instituto Nacional del Tórax, Departamento de Apoyo de Rehabilitación Cardiopulmonar Integral}{Atención de pacientes UCI, UPC, intermedio y médico-quirúrgico; evaluación instrumental de la deglución mediante FEES junto a Otorrinolaringología; evaluación de vía aérea y terapia miofuncional en apnea del sueño.}
\entry{Fonoaudiólogo}{jun. 2020--abr. 2023}{Instituto Nacional del Tórax, Servicio de Medicina Física y Rehabilitación}{Atención clínica en deglución, vía aérea y rehabilitación; evaluación instrumental; apoyo a Unidad de Medicina del Sueño; funciones de gestión de profesionales y equipos.}
\entry{Fonoaudiólogo interconsultor}{jun. 2015--jun. 2020}{Instituto Nacional del Tórax, Unidad de Medicina del Sueño}{Apoyo en evaluación instrumental de deglución y vía aérea, además de intervención miofuncional en trastornos respiratorios del sueño.}
\entry{Fonoaudiólogo - Clínica de Memoria}{ene. 2018--jun. 2020}{Hospital del Salvador, Clínica de Memoria y Neuropsiquiatría}{Evaluación e intervención en demencias atípicas y alteraciones cognitivas, ejecutivas y del lenguaje.}
\entry{Fonoaudiólogo}{mar. 2017--presente}{Laboratorio de Neuropsicología y Neurociencias Clínicas (LANNEC)}{Participación en evaluaciones cognitivas en contexto diagnóstico de demencias y afasias progresivas primarias.}
\entry{Fonoaudiólogo interconsultor}{abr. 2015--jun. 2020}{Hospital del Salvador / Instituto Nacional del Tórax, Equipo Vía Aérea y Deglución}{Atención de pacientes con trastornos de la deglución en hospitalización y apoyo en nasofibroscopía junto a Otorrinolaringología.}
\entry{Fonoaudiólogo titular}{nov. 2014--jun. 2020}{Hospital del Salvador, Servicio de Neurología, Unidad de Fonoaudiología}{Atención hospitalaria y ambulatoria de personas adultas con patología neurológica en deglución, habla y lenguaje.}
\entry{Profesional técnico}{oct. 2010--abr. 2014}{Servicio de Salud Arauco, Unidad Salud del Personal}{Gestión de salud laboral, coordinación de mesa técnica de ausentismo, gestión de listas de espera y capacitación funcionaria. Logro reportado: disminución del índice de ausentismo laboral desde 29 a 17,8 días.}
\entry{Ejercicio libre de la profesión}{2017--presente}{Santiago, Chile}{Evaluación e intervención en habla, lenguaje, deglución, voz y cognición en población infanto-juvenil y adulta.}

\section*{Experiencia académica seleccionada}
\entry{Académico invitado}{nov. 2025}{Universidad San Sebastián, Diplomado en Geriatría}{Clases sobre transición demográfica y epidemiológica, envejecimiento poblacional, gerontología y políticas públicas para personas mayores.}
\entry{Académico invitado}{jun. 2024}{Universidad Mayor, Magíster en Neuropsicología}{Clase: ``Evaluación en síndromes afásicos, una aproximación clínica''.}
\entry{Académico part-time}{ago. 2023--presente}{Universidad de los Andes, Escuela de Fonoaudiología}{Docencia en Rehabilitación Basada en la Comunidad y Salud Pública.}
\entry{Fonoaudiólogo invitado}{ago. 2022}{Universidad de Chile, Magíster en Neurociencia mención Neuropsicología}{Taller de Exploración Neuropsicológica: ``Síndromes focales del lenguaje: el caso de las afasias''.}
\entry{Docente}{ago. 2014--dic. 2018}{Universidad Internacional SEK, Carrera de Fonoaudiología}{Cátedras de Trastornos de la Deglución e Intervención Fonoaudiológica en los Trastornos de la Deglución.}
\entry{Docente clínico}{2017}{Universidad del Desarrollo / Hospital Provincial de Curanilahue}{Instructor clínico en habilitación profesional; apoyo a implementación de screening neonatal auditivo y participación en visita domiciliaria y cuidados paliativos.}
\entry{Ayudante académico}{2006--2009}{Universidad del Desarrollo}{Biofísica y Biomecánica, Física Acústica, Audiología I y apoyo a docencia de posgrado.}

\section*{Publicaciones y producción académica}
\pub{Gaete J, Ríos N, Ramírez S, \textbf{Salgado G}, et al. Six-month follow-up of the `Mi Mejor Plan' school-based prevention program: a pilot cluster randomized controlled trial among early adolescents in Chile. \emph{BMC Public Health}. 2026.}{\href{https://doi.org/10.1186/s12889-026-28548-x}{doi:10.1186/s12889-026-28548-x}}
\pub{\textbf{Salgado G}, Gaete J, Gana S, et al. Acceptability, feasibility and fidelity of the culturally adapted version of Unplugged (`Yo Sé Lo Que Quiero'): a substance use preventive program among adolescents in Chile: a pilot randomized controlled study. \emph{BMC Public Health}. 2024;24:2026.}{\href{https://doi.org/10.1186/s12889-024-19499-2}{doi:10.1186/s12889-024-19499-2}}
\pub{Gutiérrez-Arias R, \textbf{Salgado-Maldonado G}, Valdivia PL, et al. Assessing swallowing disorders in adults on high-flow nasal cannula in critical and non-critical care settings. A scoping review protocol. \emph{PLOS ONE}. 2023;18(10):e0291803.}{\href{https://doi.org/10.1371/journal.pone.0291803}{doi:10.1371/journal.pone.0291803}}
\pub{Gutiérrez-Arias R, Neculhueque-Zapata X, Valenzuela-Suazo R, et al., \textbf{Salgado-Maldonado G}, et al. Assessment of activities and participation of people by rehabilitation-focused clinical registries: a systematic scoping review. \emph{European Journal of Physical and Rehabilitation Medicine}. 2023;59(5):640--652.}{doi:10.23736/S1973-9087.23.07895-4}
\pub{Slachevsky A, Pinto A, Gajardo B, \textbf{Salgado-Maldonado G}, et al. \emph{Guía Funcionamiento Unidad de Memoria}. Centro de Memoria y Neuropsiquiatría, primera versión. 2022.}{}
\pub{Tobar-Fredes R, Briceño Meneses B, Fuentealba Miranda I, et al., \textbf{Salgado-Maldonado G}, et al. Consideraciones clínicas para fonoaudiólogos en el tratamiento de personas con COVID-19 y traqueostomía. Parte I: Deglución. \emph{Revista Chilena de Fonoaudiología}. 2020;19:1--12.}{doi:10.5354/0719-4692.2020.60185}
\pub{Tobar-Fredes R, Briceño Meneses B, Venegas-Mahn M, et al., \textbf{Salgado-Maldonado G}, et al. Consideraciones clínicas para fonoaudiólogos en el tratamiento de personas con COVID-19 y traqueostomía. Parte II: Mejorando la fonación para facilitar la comunicación. \emph{Revista Chilena de Fonoaudiología}. 2020;19:1--9.}{doi:10.5354/0719-4692.2020.60187}
\pub{Orientaciones técnicas para la rehabilitación en tiempos de pandemia COVID-19. Prevención síndrome post-COVID. Ministerio de Salud de Chile. Participación como panel revisor. 2020.}{}

\section*{Consultorías y asesorías}
\begin{itemize}
  \item Instituto Nacional del Tórax: integrante del Comité de Investigación Institucional, mar. 2025--abr. 2026.
  \item Instituto Nacional del Tórax: integrante y secretario del Comité Ético Asistencial Institucional, ene. 2025--abr. 2026.
  \item MINSAL: Panel de Expertos, GPC Rehabilitación COVID-19 - pregunta de deglución, 2022.
  \item MINSAL: Panel de Expertos, Guía GES COVID, 2021.
  \item MINSAL: participación en GPC con metodología GRADE de COVID-19 en rehabilitación, 2021.
  \item FONASA: Panel de Expertos para canasta de prestaciones de evaluación deglución, 2021.
  \item Revista Areté: revisor invitado, 2021.
\end{itemize}

\section*{Liderazgo y sociedades científicas}
\begin{itemize}
  \item \textbf{Sociedad Chilena de Alimentación y Deglución (SOCHIDA)}: socio activo desde 2020; presidente 2025--2026.
  \item \textbf{Sociedad Chilena de Medicina Intensiva (SOCHIMI)}: socio activo desde 2017; presidente de la División de Fonoaudiología y Terapia Ocupacional Intensiva, 2021--2022.
  \item \textbf{Sociedad Chilena de Medicina del Sueño (SOCHIMES)}: socio activo desde 2019.
  \item \textbf{Sociedad Chilena de Enfermedades Respiratorias (SER)}: socio activo desde 2020.
\end{itemize}

\section*{Premios y reconocimientos}
\begin{itemize}
  \item Primer lugar, trabajo oral ``Evaluación de la deglución en adultos conectados a cánula nasal de alto flujo. Una revisión de alcance'', Congreso SOCHIMI, 2024.
  \item Primer lugar, trabajo oral ``Fonoaudiología en un hospital público de alta complejidad: experiencias en tiempos de pandemia'', Congreso SOCHIMI, 2020.
\end{itemize}

\section*{Presentaciones en congresos y jornadas seleccionadas}
\begin{itemize}
  \item IV Congreso Latinoamericano de Disfagia (CLAD IV), Cancún, may. 2026. ``Algoritmos para decanulación en traqueostomía: rutas a seguir en el contexto clínico. Experiencias en el Instituto Nacional del Tórax''.
  \item Congreso de Otoño SOCHIMI, Santiago, may. 2026. ``Evidencia e innovación: optimización de la rehabilitación transdisciplinaria mediante herramientas tecnológicas en UCI''.
  \item I Curso de Voz y Deglución, Clínica Universidad de los Andes, abr. 2026. ``Deglución y enfermedad respiratoria crónica: del diagnóstico a la intervención especializada''.
  \item Congreso Anual SOCHIMI, Pucón, nov. 2025. ``Implementación de tecnologías en la rehabilitación del paciente crítico con enfoque en deglución y comunicación''.
  \item Congreso Anual SOCHIMI, Viña del Mar, oct. 2024. ``Uso de la tecnología en la rehabilitación del paciente crítico''.
  \item III Congreso Latinoamericano de Disfagia / I Congreso Chileno de Deglución y Alimentación, Santiago, jul. 2023. ``Algoritmos para decanulación en traqueostomía''.
  \item 52.º Congreso Chileno de Enfermedades Respiratorias, Santiago, nov. 2020. ``Cuidados de vía aérea superior y fonación''.
\end{itemize}

\section*{Docencia, dirección de cursos y actividades científicas seleccionadas}
\begin{itemize}
  \item Director académico, Diplomado de Deglutología del Adulto, Universidad de los Andes, desde 2026.
  \item Director académico, Diplomado de Fonoaudiología en Medicina Intensiva, Universidad de los Andes, desde 2024.
  \item Curso precongreso SOCHIMI: ``Métodos tradicionales y de bajo costo para la evaluación de la deglución, voz y vía aérea superior'', 2024.
  \item Director de programa, III Congreso Latinoamericano de Disfagia / I Congreso Chileno de Deglución y Alimentación, 2023.
  \item Director / codirector de cursos y jornadas de la División de Fonoaudiología y Terapia Ocupacional Intensiva de SOCHIMI, 2022--2023.
  \item Docente SOCHIMI: ``Elementos claves para desarrollar tu investigación'', 2021 y 2022.
\end{itemize}

\section*{Formación continua seleccionada}
\begin{itemize}
  \item Curso teórico básico de ECMO, Hospital San Juan de Dios, 40 h, 2023--2024.
  \item RCP avanzado, Instituto Nacional del Tórax, 2023.
  \item Manejo hospitalario de la persona con afasia en etapa aguda, NeuroMav, 30 h, 2023.
  \item Introducción a la ecografía pulmonar y diafragmática, Instituto Nacional del Tórax, 2022.
  \item Curso de insuficiencia respiratoria aguda y ventilación mecánica para enfermería, SOCHIMI, 55 h, 2021.
  \item Ventilación mecánica invasiva y no invasiva, Red Learning, 120 h, 2021.
  \item Certificación LSVT LOUD, 2018.
  \item Certificación en método ENMD y electroestimulación para trastornos de la deglución, 2018.
  \item Tratamiento fonoaudiológico en disturbios respiratorios del sueño: ronquido y apnea obstructiva del sueño, NICS, 36 h, 2017.
  \item Epidemiología en salud ocupacional para la red hospitalaria de salud pública, Escuela de Salud Pública Universidad de Chile, 40 h, 2013.
  \item Epidemiología en gestión de calidad y seguridad del paciente, MINSAL / TELEDUC, 100 h, 2011.
\end{itemize}

\section*{Competencias técnicas e idiomas}
\begin{tabularx}{\textwidth}{>{\bfseries}p{0.22\textwidth} X}
Herramientas & Microsoft Office avanzado; Stata I/C 14; Zotero, Mendeley, EndNote, Paperpile y Rayyan.\\
Equipamiento & Equipos audiológicos; nasofibroscopía; ecografía aplicada a vía aérea; análisis acústico de voz.\\
Idiomas & Inglés: nivel básico.\\
\end{tabularx}

\vfill
{\small\color{textgray}Versión académica pública elaborada a partir del CV actualizado proporcionado por el autor. Se omiten deliberadamente datos personales sensibles y formación escolar.}

\end{document}
